\addcontentsline{toc}{chapter}{Abstract}\vspace{-1cm}
%Border
\begin{tikzpicture}[remember picture, overlay]
  \draw[line width = 4pt] ($(current page.north west) + (0.75in,-0.75in)$) rectangle ($(current page.south east) + (-0.75in,0.75in)$);
\end{tikzpicture}


Food waste is a critical global issue, with approximately 19\% of all food produced annually being discarded before consumption. Traditional methods of food spoilage detection rely on manual inspection, which is subjective, labor-intensive, and prone to errors. Existing automated systems often use expensive electronic noses or provide only binary "fresh vs rotten" classifications without temporal context.

The objective of this work is to develop a low-cost IoT + AI system for early food spoilage detection using a combination of gas, temperature, and humidity sensors. The system employs a multi-modal approach with two main AI models: a 3-layer MLP for instant gas classification and a temporal attention-augmented LSTM for continuous shelf-life prediction in hours.

The MLP model was trained on the UCI Gas Sensor Array Drift Dataset (13,910 real sensor samples), achieving 96.84\% accuracy, 82.97\% precision, and 92.07\% recall. The LSTM model was trained on a synthetic 48-hour spoilage time-series dataset, achieving a mean absolute error of $\pm$4.0 hours as a proof-of-concept evaluation. The system also includes on-device temperature and humidity compensation for the MQ135 sensor, performed directly on the Arduino microcontroller to estimate ammonia-sensitive VOC concentration. The vision branch is deployed on the server backend using a pre-trained Vision Transformer and fine-tuned CNN.

The proposed system addresses key limitations of existing work by providing continuous shelf-life regression instead of binary classification, using a single low-cost MQ135 sensor instead of expensive sensor arrays, and incorporating temporal attention to capture spoilage dynamics over time.

\pagebreak
